% Copyright 2022  李文威 (Wen-Wei Li).
% Permission is granted to copy, distribute and/or modify this
% document under the terms of the Creative Commons
% Attribution 4.0 International (CC BY 4.0)
% http://creativecommons.org/licenses/by/4.0/

% 《代数学方法》卷一自订封面页, 由主档引入.

\setCJKfamilyfont{coverfont}{Noto Serif SC Bold}	% 设置书名字体
\setCJKfamilyfont{cover-author-font}{Noto Sans CJK SC}	% 设置作者字体
\colorlet{octa}{cyan!50!gray}	% Color for the octahedron

\begin{titlepage}\begin{tikzpicture}[remember picture, overlay, pencildraw/.style={
					color=octa, thick,
					decorate,
					decoration={random steps, segment length=1pt, amplitude=0.7pt}
				}]

		\node[anchor=center] (title) at ([xshift=19em, yshift=-10em] current page.north west)
		{ \fontsize{45}{45}\CJKfamily{coverfont}代数学方法};

		\node[anchor=center] (volume) at ([yshift=-7em] title.center) {\fontsize{30}{30}\CJKfamily{coverfont}卷一：基础架构};

		\node[anchor=west] (author) at ([yshift=-5em, xshift=4pt] volume.south west) {\fontsize{18}{18}\CJKfamily{cover-author-font}李文威};
		\node[anchor=west] at ([xshift=2.4em] author.east) {\fontsize{18}{18}\CJKfamily{cover-author-font}著};

		\draw[line width=2pt, color=black!70!gray] ([xshift=-2.5em, yshift=1em] author.north west) -- ++(25em,0);
		\shade[top color=gray, bottom color=black,] ([xshift=-1em, yshift=1em] title.north west) rectangle ++(-1em,-20em);

		% The anchors for the pictures at lower-right corner.
		\coordinate (pic) at ([xshift=-23em, yshift=22em] current page.south east);
		\coordinate (pic2) at ([xshift=3em] pic);

		% The pentagon axiom
		\begin{scope}[shift = (pic2), scale=1.5, opacity=0.55]
			\node (P0) at (90:2.3cm) {$((X \otimes Y) \otimes Z) \otimes W$};
			\node (P1) at (90+72:2cm) {$(X \otimes (Y \otimes Z)) \otimes W$} ;
			\node (P2) at (90+2*72:2cm) {\makebox[5ex][r]{$X \otimes ((Y \otimes Z) \otimes W)$}};
			\node (P3) at (90+3*72:2cm) {\makebox[5ex][l]{$X \otimes (Y \otimes (Z \otimes W))$}};
			\node (P4) at (90+4*72:2cm) {$(X \otimes Y) \otimes (Z \otimes W)$};
			\path[commutative diagrams/.cd, every arrow, every label]
			(P0) edge node[swap] {$a(X,Y,Z) \otimes \identity_W$} (P1)
			(P1) edge node[swap] {$a(X, Y \otimes Z, W)$} (P2)
			(P2) edge node[swap, inner sep=1em] {$\identity_X \otimes a(Y,Z,W)$} (P3)
			(P4) edge node {$a(X, Y, Z \otimes W)$} (P3)
			(P0) edge node {$a(X \otimes Y, Z, W)$} (P4);
		\end{scope}

		% The octahedron
		\begin{scope}[shift=(pic), rotate=-30, scale=2]
			\coordinate (A1) at (pic);
			\coordinate (A2) at (6em, 2em);
			\coordinate (A3) at (10em, 0);
			\coordinate (A4) at (4em, -2em);
			\coordinate (B1) at (5em, 5em);
			\coordinate (B2) at (5em, -5em);
		\end{scope}

		\begin{scope}[loosely dashed, opacity=0.6]
			\draw[pencildraw] (A1) -- (A2) -- (A3);
			\draw[pencildraw] (B1) -- (A2) -- (B2);
		\end{scope}
		\draw[pencildraw, fill=octa!10,opacity=0.6] (A1) -- (A4) -- (B1);
		\draw[pencildraw, fill=octa!20,opacity=0.6] (A1) -- (A4) -- (B2);
		\draw[pencildraw, fill=octa!30,opacity=0.6] (A3) -- (A4) -- (B1);
		\draw[pencildraw, fill=octa!50,opacity=0.6] (A3) -- (A4) -- (B2);
		\draw[pencildraw] (B1) -- (A1) -- (B2) --node[sloped, below] {\small\sffamily \today} (A3) --cycle;
	\end{tikzpicture}

	\clearpage	% 进入内页
	\begin{center}
		\Large{\sffamily\bfseries\thmheiti 网络版 \\ 2023 年 2 月修订} \\ \vspace{2em}
		\Large{\sffamily\bfseries\thmheiti 编译日期: \today} \\ \vspace{1em}
		%	版面: B5 (176×250mm) \\ \vspace{1em}
		本书已由高等教育出版社发行 \\
		(2019 年 1 月第 1 版, 2023 年 2 印) \\
		\texttt{ISBN: 978-7-04-050725-6}
	\end{center}
	\vfill

	\begin{flushleft} \small
		李文威 \\
		个人主页: \href{https://www.wwli.asia}{www.wwli.asia} \\
		(含勘误表等信息)
	\end{flushleft}
	\vspace{1.5em}
	\begin{tabular*}{\textwidth}{ccc}
		\includegraphics{ccby.png}
		& \begin{minipage}[b]{0.6\textwidth}
			\small\sffamily
			本作品采用知识共享 署名 4.0 国际 许可协议进行许可. 访问 \url{http://creativecommons.org/licenses/by/4.0/} 查看该许可协议.
		\end{minipage}
	\end{tabular*}

	\cleardoublepage
\end{titlepage}
